\documentclass[../../main.tex]{subfiles}
\begin{document}

\begin{flushright}
\footnotesize\textit{【自動文字起こし・要確認】}
\end{flushright}

{\bf [解]}
左から$k$両目まで見た時，$k$両目が赤である場合の数を$a_k$，青か黄色である場合の数を$b_k$とする．題意から
\begin{align}
a_{k+1}&=a_k+b_k\\
b_{k+1}&=2a_k
\end{align}
であり，又，$a_1=1$，$b_1=2$となる．※から$b_k$を消して，
\begin{align}
a_{k+2}=a_{k+1}+2a_k，\qquad a_1=1，\ a_2=3
\end{align}

だから，解いて，
\begin{align}
a_k&=\dfrac{1}{3}\left\{2^{k+1}+(-1)^k\right\}\\
b_k&=\dfrac{2}{3}\left\{2^k+(-1)^{k-1}\right\}\quad(k\ge2)
\end{align}

だから$n\ge2$の時，もとめるのは
\begin{align}
a_n+b_n&=\dfrac{1}{3}\left(2^{n+2}+(-1)^n+2(-1)^{n-1}\right)\\
&=\dfrac{1}{3}\left(2^{n+2}+(-1)^{n-1}\right)．
\end{align}

\newpage

\end{document}
